\chapter{Affine Transformations \& Hierarchical Scene Graph}

\section{The Model Matrix: Encapsulation in SceneObject}
To place any 3D entity into the world, our project implements a base class called \texttt{SceneObject} within \texttt{objects.js}. This class acts as the blueprint for everything that exists in the game, including players, goals, and the ball. Instead of scattering numbers across the code, \texttt{SceneObject} groups together three essential properties for every item: its position vector, its rotation angles along the three axes, and its scale factors.

The connection to the shaders happens through the \texttt{createModelMatrix} function. Inside this method, the code uses \texttt{gl-matrix} to build a clean $4\times4$ transformation matrix. It starts with a blank identity matrix and applies the transformations in a strict, standard order: first translating to the world position, then rotating around the X, Y, and Z axes, and finally scaling the unit geometry. This resulting model matrix is what we send directly to the WebGL vertex shader as a uniform variable. This allows the GPU to know exactly where and how to render that specific object in space.

\section{Hierarchical Structure of the Players}
One of the most complex visual challenges of the project is representing the avatars. A player is not just a single static shape [cite: 28]; it is a combination of a body, a large head, and moving feet[cite: 28, 29]. To make them move together naturally, we designed a conceptual tree structure—a local scene graph—where the torso acts as the master or "root" element[cite: 25, 28].

In this setup, the head and the feet are sub-nodes that depend directly on the body[cite: 29]. From a code perspective, when a player moves across the pitch or jumps into the air, we only update the main coordinates of the body[cite: 30]. Because of the hierarchical link, the head and limbs follow the body automatically without needing separate manual position updates for each individual part[cite: 30]. This mirrors how compound objects behave in professional graphics engines.

\section{Movement Propagation: Matrix Cloning over Traditional Stacks}
Instead of using the old-fashioned matrix stack push and pop mechanics typical of older graphics frameworks, our project uses a modern and direct approach in JavaScript by using matrix cloning via \texttt{mat4.clone()}[cite: 27].

To render the limbs solidary to the character, the drawing routine first calculates the global matrix of the body[cite: 31]. When it is time to draw a child element, like a foot that is swinging to kick the ball, the code clones the body's matrix[cite: 31]. This creates a separate snapshot of the parent's current position and rotation[cite: 27, 31]. Then, the specific local adjustments of the foot (such as its orbit position and its kicking rotation) are multiplied on top of this cloned copy[cite: 32, 33, 36]. 

The following code section from our \texttt{main.js} file shows how this branching is implemented in our \texttt{renderLeg} function to attach and animate the upper leg and the foot seamlessly relative to the hip joint:

\begin{lstlisting}[style=vscode_JS]
function renderLeg(root, hip, angle, pantsColor, shoeColor, direction) {
    // 1. Clone the parent root matrix (Body position snapshot)
    const legBase = mat4.clone(root);
    
    // 2. Translate to the hip joint location relative to the torso
    mat4.translate(legBase, legBase, hip);
    
    // 3. Apply the dynamic angular rotation for the walking animation
    mat4.rotateY(legBase, legBase, angle);

    // 4. Position and scale the Upper Leg bone structure
    const upperLeg = mat4.clone(legBase);
    mat4.translate(upperLeg, upperLeg, [0, 0, -0.15]);
    mat4.scale(upperLeg, upperLeg, [0.12, 0.12, 0.30]);
    renderObject(upperLeg, buffers.cube, pantsColor, 0.2);

    // 5. Render the Foot as a nested sub-child of the leg base matrix
    const foot = mat4.clone(legBase);
    mat4.translate(foot, foot, [direction * 0.07, 0, -0.36]);
    mat4.scale(foot, foot, [0.32, 0.16, 0.11]);
    renderObject(foot, buffers.cube, shoeColor, 0.32);
}
\end{lstlisting}

This cascading propagation ensures that the foot inherits the global movement and jump height of the torso perfectly [cite: 37], while remaining completely free to perform its own separate kicking animations without distorting the rest of the body[cite: 37].